% Part: first-order-logic
% Chapter: syntax-and-semantics
% Section: main-operator

\documentclass[../../../include/open-logic-section]{subfiles}

\begin{document}

\olfileid{fol}{syn}{mai}

\olsection{فارمولے دا \printtoken{S}{main operator}}

\begin{explain}
!!a{formula}~$!A$ بناؤن وچ سب توں آخر ورتے عملگر
دی گل کرنا اکثر فائدہ مند ہوندا اے۔ ایس عملگر نوں
~$!A$ دا \emph{مرکزی عامل} آکھدے آں۔ بدیہی طور
تے ایہ $!A$ دا ”سب توں باہرا“ عملگر اے۔ مثال
لئی، $\lnot !A$ دا مرکزی عامل $\lnot$ اے،
تے $(!A \lor !B)$ دا مرکزی عامل $\lor$ اے، وغیرہ۔
\end{explain}


\begin{defn}[!!^{main operator}]
\ollabel{def:main-op}
!!a{formula}~$!A$ دا \emph{!!{main operator}}
اینج تعریف کیتا جاندا اے:
\begin{enumerate}
\item \indcase*{!A}{!A}{$\indfrm$ دا کوئی !!{main operator} نہیں۔}

\tagitem{prvNot}{\indcase{!A}{\lnot !B}{$\indfrm$ دا !!{main operator}
  ~$\lnot$ اے۔}}{}

\tagitem{prvAnd}{\indcase{!A}{(!B \land !C)}{$\indfrm$ دا !!{main operator}
  ~$\land$ اے۔}}{}

\tagitem{prvOr}{\indcase{!A}{(!B \lor !C)}{$\indfrm$ دا !!{main operator}
  ~$\lor$ اے۔}}{}

\tagitem{prvIf}{\indcase{!A}{(!B \lif !C)}{$\indfrm$ دا !!{main operator}
  ~$\lif$ اے۔}}{}

\tagitem{prvIff}{\indcase{!A}{(!B \liff !C)}{$\indfrm$ دا !!{main operator}
  ~$\liff$ اے۔}}{}

\tagitem{prvAll}{\indcase{!A}{\lforall[x][!B]}{$\indfrm$ دا !!{main operator}
  ~$\lforall$ اے۔}}{}

\tagitem{prvEx}{\indcase{!A}{\lexists[x][!B]}{$\indfrm$ دا !!{main operator}
  ~$\lexists$ اے۔}}{}
\end{enumerate}
\end{defn}

ہر صورت وچ ساڈی مراد فارمولے دے !!{main operator}
دی دکھائی گئی خاص \emph{آمد} اے۔ مثال لئی، کیوں جے
فارمولا $((!D \lif !E) \lif (!E \lif !D))$، $(!B \lif !C)$
دے روپ وچ اے، جتھے $!B$، $(!D \lif !E)$ اے تے
$!C$، $(!E \lif !D)$ اے، ایس لئی $\lif$ دی
دوجی آمد !!{main operator} اے۔

\begin{explain}
% PNB-SYN-007: اگلے حصے دے صریح فرق مطابق ایتھے چھوٹے فارمولیاں اُتے ایس فنکشن دیاں قدراں نہیں ورتیاں؛ ایہ صورتاں دے لحاظ نال سیدھی تعریف اے، بازگشتی تعریف نہیں۔
ایہ اک فنکشن دی \emph{صورتاں دے لحاظ نال سیدھی} تعریف اے، جیہڑا سارے
غیر ایٹمی !!{formula}s نوں اوہناں دے !!{main operator}
دی آمد نال جوڑدا اے۔ جس طرح !!{formula}s
استقرائی طور تے تعریف کیتے گئے نیں، ہر
!!{formula}~$!A$، \olref{def:main-op} دیاں
صورتاں وچوں اک پوری کردا اے۔ ایہ ضمانت دیندا
اے کہ ہر غیر ایٹمی !!{formula}~$!A$ لئی
!!a{main
  operator} موجود اے۔ کیوں جے ہر !!{formula}
ایہناں شرطاں وچوں صرف اک پوری کردا اے، تے
ہر صورت وچ اوہ چھوٹے !!{formula}s وی اکّو طرح
متعین نیں جیہناں توں $!A$ بنیا اے، اس لئی
~$!A$ دے !!{main
  operator} دی آمد اکّو اے، تے ایس طرح
اسیں اک فنکشن تعریف کر دتا اے۔
\end{explain}

اسیں !!{formula}s نوں \olref{tab:main-op} وچ
دتے ناں اوہناں دے !!{main operator} دی علامت
مطابق دیندے آں۔\iftag{defNot,defOr,defAnd,defIf,defIff,defTrue,defFalse,defEx,defAll}
{ پر یاد رکھو کہ تعریف شدہ عملگر باقاعدہ طور تے
!!{formula}s وچ نہیں آؤندے۔ اوہ صرف مخفف نیں،
اس لئی باقاعدہ طور تے کسے فارمولے دے مرکزی عامل
نہیں بن سکدے۔ لہٰذا سارے !!{formula}s بارے ثبوتاں
وچ اوہناں دا وکھرا معاملہ کرن دی لوڑ نہیں۔}

% PNB-SYN-002 ماخذ دی درستی: جدول دے عطف، فصل تے شرطیہ مثالاں دیاں بند ہُندیاں قوساں ماخذ وچ ریاضی موڈ توں باہر نیں؛ ایتھے تیناں مکمل فارمولے ریاضی موڈ وچ رکھے نیں۔
\begin{table}[!h]
\centering
\begin{tabular}{c | c | c}
!!^{main operator} & !!{formula} دی قسم & مثال\\
\hline
کوئی نہیں & ایٹمی (!!{formula}) &
\iftag{prvFalse}{$\lfalse$,}{}
\iftag{prvTrue}{$\ltrue$,}{}
$\Atom{R}{t_1, \dots, t_n}$,
$\eq[t_1][t_2]$\\
$\lnot$ & نفی & $\lnot !A$ \\
$\land$ & عطف & $(!A \land !B)$ \\
$\lor$ & فصل & $(!A \lor !B)$ \\
$\lif$ & !!{conditional} & $(!A \lif !B)$ \\
$\liff$ & !!{biconditional} & $(!A \liff !B)$ \\
$\lforall[][]$ & کُلّی (!!{formula}) & $\lforall[x][!A]$ \\
$\lexists[][]$ & وجودی (!!{formula}) & $\lexists[x][!A]$
\end{tabular}
\caption{!!{formula}s دے ناں تے مرکزی عامل}
\ollabel{tab:main-op}
\end{table}

\end{document}
